% Management in Global Markets – Midterm Cheat Sheet
% Compile with: pdflatex midterm-cheat-sheet.tex

\documentclass[9pt,landscape,a4paper]{article}
\usepackage[utf8]{inputenc}
\usepackage[T1]{fontenc}
\usepackage[left=1cm,right=2.8cm,top=1cm,bottom=1cm,marginparwidth=2.1cm,marginparsep=0.35em]{geometry}
\usepackage{multicol}
\usepackage{marginnote}
\usepackage{enumitem}
\usepackage{parskip}
\usepackage{xcolor}
\usepackage{soul}
\usepackage{fancyhdr}
\definecolor{lightpink}{RGB}{255,182,193}
\definecolor{lightpeach}{RGB}{255,218,185}
\definecolor{noteblue}{RGB}{173,216,230}
\newcommand{\handnote}[1]{\marginnote{\raggedright\small\color{noteblue}#1\vspace{0.6em}}}
\newcommand{\hlpink}[1]{\sethlcolor{lightpink}\hl{#1}\sethlcolor{lightpeach}}
\newcommand{\hlpeach}[1]{\sethlcolor{lightpeach}\hl{#1}}
\sethlcolor{lightpeach}
\setlength{\parskip}{0.3em}
\setlength{\columnsep}{0.8em}
\pagestyle{fancy}
\fancyhf{}
\fancyfoot[R]{\raisebox{0.8em}{\thepage}}
\fancypagestyle{plain}{\fancyhf{}\fancyfoot[R]{\raisebox{0.8em}{\thepage}}}
\raggedright

\begin{document}
\begin{center}
{\large\bfseries Management in Global Markets -- Midterm Cheat Sheet}\\[0.3em]
\small\textit{Define terms, give one example, tie to internationalization.}\\[0.6em]
\fbox{\parbox{0.92\textwidth}{\centering\small
\textcolor{black!80}{\textbf{Key:}} \colorbox{lightpink}{\strut\ \ \ } \textbf{Pink} = definitions \quad
\colorbox{lightpeach}{\strut\ \ \ } \textbf{Peach} = frameworks \quad
\colorbox{noteblue}{\strut\ \ \ } \textbf{Blue} = margin notes
\par}}
\end{center}
\vspace{-0.4em}

\begin{multicols}{2}

\section*{Unit 1 -- Globalization}\handnote{Winners vs.\ losers.}
\begin{itemize}[leftmargin=*,nosep]
\item \textbf{\hlpink{Globalization}} = Free movement. High production + export.
\item \textbf{\hlpink{Slowbalization}} = Slowing. Protectionist. Localization.
\item Winners: China, India, Indonesia. USA = biggest importer.
\item \textbf{\hlpink{Economic complexity}} = Raw $\to$ complex products. Tech, R\&D. USA, UK, Netherlands.
\item Demographics: Europe ageing. Africa growing. Young = production.
\end{itemize}

\section*{Unit 2 -- Supply Chain}\handnote{Quality vs.\ cost.}
\begin{itemize}[leftmargin=*,nosep]
\item \textbf{\hlpink{Supply chain}} = Parts move before product made. Global.
\item Japan = quality. China = cost, volume.
\item Risks: culture, laws, language. Competition: go where competitors are.
\item Culture shock: excitement $\to$ hard $\to$ adapt.
\end{itemize}

\section*{Unit 3 -- Strategy Levels}\handnote{Corporate, competitive, functional.}
\begin{itemize}[leftmargin=*,nosep]
\item \textbf{Corporate}: Which businesses? (Cars vs.\ trucks).
\item \textbf{Competitive}: How to compete? (Price, quality, niche).
\item \textbf{Functional}: HR, marketing, finance, ops, R\&D.
\item Keep simple. Long-term. 15 years, not 15 days. ``Slow is fast.''
\end{itemize}

\section*{Unit 4 -- Internationalization}\handnote{Location = critical.}
\begin{itemize}[leftmargin=*,nosep]
\item \textbf{\hlpink{SME}} = $<$250 employees, $<$50M. Spain: 99\% SMEs.
\item \textbf{Location choice}: Critical. Wrong = ruin. Politics, economy, market.
\item \textbf{Resources} = what you have. \textbf{Capabilities} = what you do.
\item \textbf{Competitive advantage} = better than rivals. \textbf{Absolute} = almost no one has. Rare.
\item Before going abroad: strong CA? If not, stay home.
\item Entry: acquisition, JV, alliances (licensing, distribution).
\end{itemize}

\section*{Unit 5 -- Chile Case}\handnote{Tech transfer risk.}
\begin{itemize}[leftmargin=*,nosep]
\item French company $\to$ Chile. Lost money, jobs, technology.
\item Local partners used tech independently. No protection. Communication broke.
\item Lesson: careful with core tech. IP, contracts, partners. Culture, trust.
\end{itemize}

\section*{Unit 6 -- Export vs.\ Int'l \& HRM}\handnote{People matter.}
\begin{itemize}[leftmargin=*,nosep]
\item \textbf{\hlpink{Export}} = Product moves. Company home. Low risk.
\item \textbf{\hlpink{Internationalization}} = Operations abroad. Plants, subsidiaries. FDI.
\item \textbf{FDI} = Long-term commitment. Needs international HRM.
\item HRM: attract, select, train, move talent. Trust, promotion.
\item Motivation: salary + ambition, recognition. As important as skills.
\item Who goes? Motivated, resilient. Ready for sacrifice.
\end{itemize}

\section*{Quick revision}\handnote{Run through before exam.}
\begin{itemize}[leftmargin=*,nosep]
\item Global vs.\ slowbalization. Supply chain. Strategy levels.
\item SME. Location. Resources vs.\ capabilities. CA vs.\ absolute.
\item Chile. Export vs.\ int'l. FDI. HRM. Motivation.
\end{itemize}

\section*{\hlpeach{Common question angles}}\handnote{Memorise these.}
\begin{itemize}[leftmargin=*,nosep]
\item \textbf{``Why might int'l fail?''} Wrong location, weak CA, tech transfer, culture. Chile.
\item \textbf{``Strategy levels?''} Corporate / competitive / functional.
\item \textbf{``Export vs.\ int'l?''} Product moves vs.\ company moves.
\item \textbf{``Success abroad?''} Strong CA, location, people, careful tech, partners.
\end{itemize}

\end{multicols}

\newpage
\vspace*{0.3em}
\noindent\color{noteblue}
\setlength{\parskip}{0.4em}
\begin{multicols}{2}

\textbf{Globalization}\\
Global vs.\ slowbalization. Economic complexity.

\textbf{Supply chain}\\
Quality vs.\ cost. Culture. Competition.

\textbf{Strategy}\\
Corporate. Competitive. Functional.

\textbf{Internationalization}\\
SME. Location. CA. Chile case.

\textbf{Export vs.\ Int'l}\\
Product vs.\ company. FDI. HRM.

\end{multicols}
\vfill
\normalsize\normalfont\color{black}
\end{document}
